\subsection{Fundamental Theorem of Algebra} 
\begin{theorem}[Fundamental Theorem of Algebra]
For any polynomial
\[
\sum_{k=0}^n{a_n x^n}
\]
with \emph{complex} coefficients, there is a unique list of complex numbers $z_1,z_2,...,z_n$ such that
\[
\sum_{k=0}^n{a_n x^n}
=
a_n(x-z_1)(x-z_2)\cdots(x-z_n)
\]
\end{theorem}
We're not going to prove this theorem---it takes a \emph{lot} of tools we don't have yet!---but we do need to understand what it means.
\begin{itemize}
\item A polynomial with degree $n$ has exactly \makeblanksmall roots (``up to multiplicity'', which means that \makeblank[3in])
\item If we know all the roots and the leading coefficient, they ``uniquely determine'' the polynomial, which means \makeblank[3in]
\item Like numbers, polynomials factor \emph{uniquely}. That means if we want to factor a polynomial, it doesn't matter \makeblank.
\end{itemize}
\begin{Example}
How many roots does the polynomial $x^5-7x^4+3x^2+2$ have?
\begin{lpic}{}{2}
\end{lpic}
\end{Example}
We can put this together with \makeblank[3in] to make some guesses about how many of a polynomial's roots might be imaginary (that is, \makeblank[1in] numbers that are not \makeblank[1in]).
\begin{Example}
Analyze the roots of the polynomial $2x^6-7x^3-x^2-4$ have?
\begin{itemize}
\item This polynomial has \makeblanksmall sign change, so it has \makeblanksmall positive real root.
\item Substituting $-x$ give us \makeblank
\item This has \makeblanksmall sign change, so the orginal polynomial has \makeblanksmall negative real root.
\item The polynomial has degree \makeblanksmall, so it has \makeblanksmall total roots.
\item Therefore it has \makeblank imaginary roots.
\end{itemize}
\end{Example}